\documentclass[10pt,a4paper]{article}

\usepackage[a4paper,top=17mm,bottom=18mm,left=16mm,right=16mm,headheight=14pt]{geometry}
\usepackage{fontspec}
\usepackage{xcolor}
\usepackage{array}
\usepackage{tabularx}
\usepackage{enumitem}
\usepackage{fancyhdr}
\usepackage{hyperref}
\usepackage{needspace}

\setmainfont{Arial}
\setsansfont{Arial}

\definecolor{Primary}{HTML}{235347}
\definecolor{Accent}{HTML}{E98B2A}
\definecolor{Soft}{HTML}{EEF4F1}
\definecolor{LineGray}{HTML}{B7C3BE}
\definecolor{TextGray}{HTML}{4E5A56}

\newcommand{\documentlabel}{Pedoman Wawancara Intern Business Consultant Nasi Gerilya}

\hypersetup{
  colorlinks=true,
  linkcolor=Primary,
  urlcolor=Primary,
  pdftitle={Pedoman Wawancara Intern Business Consultant Nasi Gerilya},
  pdfauthor={Instrumen Penelitian Skripsi}
}

\pagestyle{fancy}
\fancyhf{}
\fancyhead[L]{\small\textcolor{Primary}{\documentlabel}}
\fancyhead[R]{\small\textcolor{TextGray}{Versi Juli 2026}}
\fancyfoot[C]{\small\textcolor{TextGray}{Halaman \thepage}}
\renewcommand{\headrulewidth}{0.4pt}
\renewcommand{\footrulewidth}{0pt}

\setlength{\parindent}{0pt}
\setlength{\parskip}{4pt}
\renewcommand{\arraystretch}{1.18}
\hyphenpenalty=10000
\exhyphenpenalty=10000
\emergencystretch=3em
\sloppy
\newcolumntype{Y}{>{\raggedright\arraybackslash}X}

\newcommand{\sectionbar}[1]{%
  \Needspace{4\baselineskip}
  \vspace{4mm}
  \noindent\colorbox{Soft}{%
    \parbox{\dimexpr\linewidth-2\fboxsep\relax}{%
      \large\bfseries\textcolor{Primary}{#1}
    }%
  }\par\vspace{2mm}
}

\newcommand{\probe}[1]{%
  \vspace{1mm}
  \noindent{\small\textbf{\textcolor{Accent}{Probing:}} \textcolor{TextGray}{#1}}\par
}

\newcommand{\answerline}{%
  \par\vspace{0.8mm}
  \noindent\textcolor{LineGray}{\rule{\linewidth}{0.35pt}}\par\vspace{2.8mm}
}

\newcommand{\questionblock}[4]{%
  \Needspace{8\baselineskip}
  \noindent\textbf{\textcolor{Primary}{#1. #2}}\par
  \noindent #3\par
  \probe{#4}
  \answerline\answerline\answerline
}

\newcommand{\metadatafield}[1]{%
  \textbf{\textcolor{Primary}{#1}}\par
  \textcolor{LineGray}{\rule{\linewidth}{0.35pt}}\par\vspace{2.4mm}
}

\begin{document}

\begin{titlepage}
\thispagestyle{empty}
\vspace*{9mm}
{\Huge\bfseries\textcolor{Primary}{Pedoman Wawancara}}\par
{\Huge\bfseries\textcolor{Primary}{Intern / Magang}}\par
{\Huge\bfseries\textcolor{Primary}{Business Consultant}}\par
\vspace{3mm}
{\Large\textcolor{Accent}{Nasi Gerilya}}\par
\vspace{7mm}
\textcolor{Accent}{\rule{0.92\linewidth}{1.4pt}}\par
\vspace{8mm}

\begin{tabularx}{0.92\linewidth}{>{\bfseries\color{Primary}}p{35mm}Y}
Jenis instrumen & Pedoman wawancara semi-terstruktur \\
Target informan & Intern atau magang \textit{business consultant} yang pernah mengamati, mendampingi, atau memberi rekomendasi untuk Nasi Gerilya \\
Fokus penelitian & Strategi pemasaran Nasi Gerilya pada GrabFood untuk meningkatkan penjualan berdasarkan analisis SWOT \\
Tujuan wawancara & Menggali diagnosis bisnis, rekomendasi pemasaran, kendala implementasi, indikator evaluasi, serta masukan IFAS dan EFAS \\
\end{tabularx}

\vfill
{\large\bfseries\textcolor{Primary}{Instrumen Penelitian Skripsi}}\par
{\normalsize\textcolor{TextGray}{Disusun untuk kebutuhan pengumpulan data lapangan, 2026}}\par
\end{titlepage}

\sectionbar{Identitas Wawancara}

\begin{tabularx}{\linewidth}{X X}
\metadatafield{Nama informan} & \metadatafield{Tanggal wawancara} \\
\metadatafield{Posisi / peran} & \metadatafield{Durasi keterlibatan} \\
\metadatafield{Tempat / media wawancara} & \metadatafield{Kode informan} \\
\end{tabularx}

\sectionbar{Catatan untuk Pewawancara}

Gunakan pedoman ini secara semi-terstruktur. Pertanyaan utama dapat ditanyakan berurutan, tetapi peneliti boleh menyesuaikan alur sesuai jawaban informan. Probing dipakai bila jawaban masih terlalu umum, belum menyebut contoh, atau belum menjelaskan dasar data.

\begin{itemize}[leftmargin=*,itemsep=2pt,topsep=2pt]
  \item Minta informan menjelaskan pengalaman aktual selama magang, bukan hanya pendapat umum.
  \item Pisahkan temuan berbasis data dengan asumsi atau saran pribadi.
  \item Catat contoh konkret, angka, dokumen, atau observasi yang disebut informan.
  \item Hubungkan jawaban dengan kekuatan, kelemahan, peluang, dan ancaman Nasi Gerilya.
\end{itemize}

\sectionbar{Pemetaan Tema Wawancara}

\begin{tabularx}{\linewidth}{|p{31mm}|Y|Y|}
\hline
\textbf{\textcolor{Primary}{Tema}} & \textbf{\textcolor{Primary}{Fokus Data}} & \textbf{\textcolor{Primary}{Kaitan Analisis}} \\ \hline
Profil peran & Ruang lingkup magang, tugas, koordinasi, dan batas keterlibatan informan. & Menentukan kredibilitas dan batas informasi informan. \\ \hline
Diagnosis bisnis & Data penjualan, pelanggan, menu, ulasan, promo, operasional, dan kompetitor. & Menjadi dasar pembacaan masalah pemasaran. \\ \hline
Strategi GrabFood & Segmentasi, positioning, produk, harga, promo, tampilan toko digital, dan proses layanan. & Menggali strategi pemasaran pada platform GrabFood. \\ \hline
Persaingan & Perbandingan dengan rumah makan Padang atau kompetitor makanan sejenis di sekitar lokasi dan platform. & Menjadi input peluang dan ancaman eksternal. \\ \hline
Evaluasi dan SWOT & KPI, hambatan implementasi, kekuatan, kelemahan, peluang, ancaman, dan prioritas rekomendasi. & Menjadi bahan IFAS, EFAS, dan matriks SWOT. \\ \hline
\end{tabularx}

\sectionbar{A. Profil Peran dan Konteks Magang}

\questionblock{W-BC-01}{Peran informan}{Bisa dijelaskan peran Anda sebagai intern atau magang \textit{business consultant} di Nasi Gerilya?}{Sejak kapan terlibat, tugas utama apa yang dikerjakan, dan dengan siapa paling sering berkoordinasi?}

\questionblock{W-BC-02}{Tujuan pendampingan}{Masalah bisnis atau pemasaran apa yang menjadi fokus utama selama Anda membantu Nasi Gerilya?}{Apakah fokusnya peningkatan penjualan, promosi, menu, operasional, kanal online, kompetitor, atau hal lain?}

\questionblock{W-BC-03}{Batas keterlibatan}{Bagian mana dari kegiatan Nasi Gerilya yang Anda amati langsung, dan bagian mana yang hanya Anda ketahui dari cerita atau data pihak lain?}{Apakah Anda ikut observasi toko, membaca data penjualan, berdiskusi dengan pemilik, atau memantau platform GrabFood?}

\sectionbar{B. Diagnosis Bisnis dan Data Awal}

\questionblock{W-BC-04}{Sumber data}{Data atau informasi apa saja yang Anda gunakan untuk memahami kondisi Nasi Gerilya?}{Data pesanan, omzet, menu terlaris, ulasan, rating, promo, observasi lapangan, wawancara internal, atau kompetitor.}

\questionblock{W-BC-05}{Masalah utama}{Dari data atau pengamatan tersebut, masalah apa yang paling terlihat pada pemasaran Nasi Gerilya?}{Apakah masalahnya di visibilitas toko, harga, promo, menu, stok, kecepatan layanan, rating, atau persaingan?}

\questionblock{W-BC-06}{Peran GrabFood}{Bagaimana penilaian Anda terhadap peran GrabFood dalam penjualan Nasi Gerilya?}{Apa kelebihan dan kekurangan GrabFood dibanding pembelian langsung atau platform pesan antar lain?}

\sectionbar{C. Strategi Pemasaran GrabFood}

\questionblock{W-BC-07}{Pelanggan target}{Menurut Anda, pelanggan seperti apa yang paling potensial disasar Nasi Gerilya melalui GrabFood?}{Pekerja sekitar, mahasiswa, keluarga, pelanggan promo, pelanggan sekitar lokasi, pelanggan makan siang, atau segmen lain.}

\questionblock{W-BC-08}{Positioning}{Citra apa yang sebaiknya paling dikuatkan oleh Nasi Gerilya di GrabFood?}{Porsi besar, rasa lauk, masakan Padang praktis, harga promo, jarak dekat, menu tertentu, atau kecepatan layanan.}

\questionblock{W-BC-09}{Menu unggulan}{Menu apa yang paling layak dijadikan andalan untuk menarik pesanan GrabFood?}{Dendeng bakar, ayam pop, ayam goreng, paket nasi, add-on, menu hemat, atau variasi menu lain.}

\questionblock{W-BC-10}{Harga dan promo}{Bagaimana penilaian Anda terhadap harga, diskon, voucher, dan ongkir Nasi Gerilya di GrabFood?}{Apakah harga online terasa kompetitif, promo mendorong pesanan, dan margin masih aman?}

\questionblock{W-BC-11}{Tampilan toko digital}{Bagian apa dari tampilan toko Nasi Gerilya di GrabFood yang perlu diperbaiki?}{Foto menu, nama menu, deskripsi, kategori, status stok, rating, ulasan, atau susunan menu.}

\questionblock{W-BC-12}{Aktivitas promosi}{Promosi seperti apa yang menurut Anda paling cocok untuk mendorong pesanan GrabFood Nasi Gerilya?}{Voucher, bundling, paket jam makan siang, konten media sosial, menu hero, atau kerja sama promosi.}

\sectionbar{D. Persaingan, Layanan, dan Operasional}

\questionblock{W-BC-13}{Kompetitor relevan}{Kompetitor mana yang paling relevan dibandingkan dengan Nasi Gerilya, dan apa alasannya?}{Bandingkan dari rating, jumlah ulasan, harga, jarak, menu unggulan, reputasi, dan promo.}

\questionblock{W-BC-14}{Pembeda usaha}{Apa pembeda utama Nasi Gerilya yang bisa dijadikan kekuatan strategi?}{Apakah pembeda berasal dari rasa, porsi, lokasi, menu tertentu, layanan, pelanggan loyal, atau fleksibilitas promosi?}

\questionblock{W-BC-15}{Kendala layanan}{Kendala operasional apa yang paling mungkin memengaruhi kepuasan pelanggan GrabFood?}{Waktu proses, stok menu, item kurang, catatan pelanggan, kemasan, koordinasi admin-dapur, atau penanganan komplain.}

\sectionbar{E. Evaluasi, SWOT, dan Rekomendasi}

\questionblock{W-BC-16}{Kelemahan utama}{Kelemahan apa yang paling perlu diperbaiki agar strategi GrabFood lebih efektif?}{Konsistensi menu, harga, tampilan digital, stok, kecepatan layanan, data, SDM, atau pengelolaan promo.}

\questionblock{W-BC-17}{Peluang dan ancaman}{Peluang apa yang masih bisa dimanfaatkan, dan ancaman apa yang paling perlu diantisipasi Nasi Gerilya?}{Peluang dari pelanggan sekitar, jam makan siang, promo platform, ulasan positif; ancaman dari pesaing, biaya platform, harga bahan, rating negatif, atau pelanggan sensitif harga.}

\questionblock{W-BC-18}{Indikator evaluasi}{Indikator apa yang sebaiknya dipakai untuk menilai keberhasilan strategi GrabFood?}{Jumlah pesanan, omzet, rata-rata nilai transaksi, repeat order, rating, komplain, menu habis, biaya promo, atau margin.}

\questionblock{W-BC-19}{Prioritas strategi}{Jika hanya boleh memilih tiga prioritas perbaikan, apa yang harus dilakukan Nasi Gerilya terlebih dahulu?}{Mengapa prioritas itu dipilih, seberapa mudah dijalankan, dan apa dampaknya terhadap penjualan atau kepuasan pelanggan?}

\questionblock{W-BC-20}{Validasi rekomendasi}{Dari semua temuan dan rekomendasi Anda, mana yang paling kuat datanya dan mana yang masih perlu diuji lagi?}{Data tambahan apa yang perlu dikumpulkan agar rekomendasi lebih meyakinkan?}

\sectionbar{Pertanyaan Penutup}

\questionblock{W-BC-21}{Tambahan informan}{Apakah ada hal penting tentang strategi Nasi Gerilya di GrabFood yang belum sempat ditanyakan?}{Minta informan menambahkan catatan bebas, contoh kasus, atau saran praktis yang dianggap penting.}

\end{document}
